\section{Two Letters and a Gospel}

The New Testament carries two letters under Peter's name, and a third
document, the Gospel of Mark, is attached to him by an ancient and unbroken
tradition. Their evidential standing must be graded separately, not
presumed, and the grades are genuinely different.

\subsection{The First Epistle}

The First Epistle of Peter addresses ``the strangers dispersed through
Pontus, Galatia, Cappadocia, Asia, and Bithynia'' (1~Pt 1:1)---five Roman
provinces covering most of Asia Minor north and west of the Taurus. It is
written to communities under social pressure rather than formal
persecution: the readers face ``the burning heat which is to try you''
(4:12) and the possibility of suffering ``as a Christian'' for the name
itself (4:16), which the letter treats as an occasion for glorifying God
rather than as a legal emergency.

The writer speaks as ``myself also an ancient, and a witness of the
sufferings of Christ'' (5:1)---in Greek \textit{sympresbyteros}, a
fellow-elder, and \textit{martys} of Christ's sufferings---and exhorts the
elders to ``feed the flock of God,'' \textit{poimanate to poimnion tou
theou}, ``taking care of it'' (5:2), where the participle is
\textit{episkopountes}, exercising oversight. Both verbs are the vocabulary
of John 21 and of the developing church order at once. He names Silvanus at
the close---``By Silvanus, a faithful brother unto you, as I think, I have
written briefly'' (5:12)---and signs off: ``The church that is in Babylon,
elected together with you, saluteth you: and so doth my son Mark'' (5:13).

Three questions have to be kept apart.

\textbf{Where was it written?} From the second century onward ``Babylon'' was
read as a figure for Rome. Eusebius reports the reading from Clement of
Alexandria and Papias: ``they say that he wrote it in Rome itself, as is
indicated by him, when he calls the city, by a figure, Babylon''
(\work{Church History} 2.15.2). The figure is not arbitrary---Babylon was the
power that destroyed the first temple, and Rome had now destroyed the
second---but the reading is an inference, and it is a second-century
inference about a first-century text, not an independent report. It is
corroborating rather than probative.

\textbf{Who put the Greek together?} The letter's Greek is competent,
periodic, and steeped in the Septuagint, which is the principal reason much
modern criticism doubts direct composition by a Galilean fisherman. Two
cautions belong here. First, the strongest textual basis for the ``fisherman
could not have written this'' argument is Acts 4:13, and, as shown above,
\textit{agrammatoi kai idiotai} on the lips of hostile officials means
untrained rather than illiterate; it is a slender foundation for a
literary judgment. Second, the Silvanus clause cannot settle the question in
either direction. The Greek \textit{dia Silouanou\ldots\ egrapsa} is the
standard formula for identifying the bearer of a letter, and it is at least
as naturally read that way as of a secretary given free compositional rein.
The honest position is that direct composition, dictated or secretarial
composition under Peter's authority, and posthumous composition within a
Petrine circle at Rome all remain defensible, and that no argument in this
study turns on the choice.

\textbf{How did antiquity receive it?} Without hesitation. ``One epistle of
Peter, that called the first, is acknowledged as genuine,'' Eusebius says,
``and this the ancient elders used freely in their own writings as an
undisputed work'' (\work{Church History} 3.3.1). This study cites the letter
as canonical Petrine teaching whose exact literary authorship is unresolved,
and uses it for two historical purposes only: as evidence that a Roman
association for Peter existed early, and as evidence of how the pastoral
charge of John 21 was understood by those writing in his name.

\subsection{The Second Epistle}

The Second Epistle is a different case, and the difference is not a modern
discovery.

It is framed as the apostle's testament---``as long as I am in this
tabernacle, to stir you up by putting you in remembrance. Being assured that
the laying away of this my tabernacle is at hand, according as our Lord Jesus
Christ also hath signified to me'' (1:13--14), an evident echo of the
prophecy in Jn 21:18--19. It claims Transfiguration eyewitness in so many
words: ``we were eyewitnesses of his greatness\ldots\ this voice we heard
brought from heaven, when we were with him in the holy mount'' (1:16--18). It
addresses a problem that presupposes the passing of the first generation---``since the time that the fathers slept, all things continue as they were
from the beginning'' (3:4)---and it treats Paul's letters as a collection
already being read alongside ``the other scriptures'' (3:15--16).

To these has to be added a literary fact that a reader can check without any
scholarly apparatus, and that this study checked directly in the Greek. The
second chapter of 2~Peter and the letter of Jude share not merely themes but
sequences and wording. Both invoke the angels who sinned,
consigned in chains of darkness and kept for judgment (2~Pt 2:4; Jude 6);
both then invoke Sodom and Gomorrha as an example (2~Pt 2:6; Jude 7); and the
description of the false teachers ends, in both, with the same Greek
clause---\textit{hois ho zophos tou skotous teteretai}, ``to whom the mist of
darkness is reserved'' (2~Pt 2:17; Jude 13), which the Byzantine text extends
in both places with \textit{eis aiona}, for ever, and the Nestle--Aland text
in Jude only. In one case the clause is attached to waterless springs and
storm-driven clouds, in the other to waterless clouds and wandering stars.
Two short documents do not share a whole clause by accident.
One used the other, or both used a common source; the majority critical
judgment is that 2~Peter used Jude, and the direction of dependence is not
something this study can establish independently.

Ancient doubt is therefore not an invention of modern criticism. Eusebius
classed the letter among the disputed books, and Jerome, cataloguing Peter's
works in 392, records that the second epistle ``on account of its difference
from the first in style, is considered by many not to be by him''
(\work{On Illustrious Men} 1). A further and more bounded negative deserves
notice: the Muratorian Fragment, the oldest surviving list of New Testament
books, generally dated to about 170 and Roman in origin, names the epistle of
Jude, two epistles of John, the Wisdom of Solomon, and the apocalypses of
John and of Peter---and does not name either epistle of Peter. The fragment
is damaged at both ends and its silence cannot be pressed; but it is the
earliest canon list we have, it comes from the very church that claimed Peter
as its own, and it does not list his letters.

The counterweight is honest too. The letter opens with the Semitic spelling
\textit{Symeon} Peter (1:1), a form applied to the apostle only once
elsewhere in the New Testament, on James's lips at the Jerusalem council
(Acts 15:14). A forger reaching for verisimilitude might have chosen it; so
might a writer standing close to the tradition in which it was natural. The
datum cuts both ways and is reported here as such.

Most modern critical scholarship regards 2~Peter as the latest writing in the
New Testament, composed in Peter's name in the tradition of the testamentary
genre; a minority defends direct authenticity. Its canonical authority is a
judgment of the Church about the book and stands on its own basis. This study
does not use 2~Peter as uncontested autobiography, and in particular does not
use its Transfiguration claim as an independent eyewitness attestation of
that scene.

\subsection{The Gospel of Mark}

One more document belongs to the Petrine dossier by tradition. The earliest
statement is quoted by Eusebius from Papias of Hierapolis, who attributes it
to ``the presbyter'':

\begin{studyframe}\small
``Mark, having become the interpreter of Peter, wrote down accurately,
though not in order, whatsoever he remembered of the things said or done by
Christ. For he neither heard the Lord nor followed him, but afterward, as I
said, he followed Peter, who adapted his teaching to the needs of his
hearers.'' (Papias, in Eusebius, \work{Church History} 3.39.15)
\end{studyframe}

Irenaeus, about 180, agrees: ``Mark, the disciple and interpreter of Peter,
did also hand down to us in writing what had been preached by Peter''
(\work{Against Heresies} 3.1.1). The Muratorian Fragment's opening lines,
which survive only in mutilated form, appear to say something similar about
Mark's presence at Peter's teaching.

What is established and what is not should be stated carefully. Established:
that by the first quarter of the second century, in Asia Minor, the second
Gospel was already connected to Peter's preaching through Mark, and that
Irenaeus at Lyons repeats the connection independently of Papias'
formulation. Also established: that the same Mark appears beside Peter in
1~Pt 5:13 at ``Babylon,'' so that the tradition and the letter agree on the
association whatever one concludes about either. Not established: that the
second Gospel is in fact Petrine memoir at one remove. Papias is writing at
two removes---he reports what ``the presbyter'' said---and his defensive
phrase ``though not in order'' shows that he is already answering a criticism
of the book rather than describing it neutrally. Whether the Gospel's content
reflects Peter's preaching is among the most debated questions in Gospel
criticism, and this study takes no position on it.

\begin{studybox}{Project synthesis --- the letters and the historical Peter}
\textbf{Judgment.} 1~Peter is usable historical evidence for the way the
Petrine tradition understood itself within a generation of the apostle's
death, and for the early association of Peter with Rome; it is not usable as
a transcript of the apostle's own Greek. 2~Peter is usable evidence for
second-generation reception of Peter's memory and for the early collection of
Paul's letters; it is not usable as autobiography at all.

\textbf{Reasons.} (1) Antiquity received the two letters differently, and the
distinction runs from Eusebius through Jerome; a study that flattens them is
less careful than the fourth century was. (2) 2~Peter's shared clause with
Jude establishes a literary relationship with another New Testament document,
which no theory of direct apostolic composition explains comfortably in
either direction. (3) The arguments about 1~Peter's Greek are real but
weaker than usually presented, since their strongest textual prop is a
hostile aside in Acts and their standard solution, the Silvanus clause, will
not carry the weight placed on it.

\textbf{Strongest counterargument.} That this grading imports a modern
literary standard into an ancient practice. Ancient authors used secretaries
with wide latitude; a letter dictated in substance and written up by a
competent hand was that author's letter by every convention of the time, and
the pseudonymity thesis assumes a modern notion of authorship that the first
century did not hold. On that view 1~Peter is simply Peter's letter, and even
2~Peter may be the work of a disciple executing an authorized commission
rather than a fiction. The argument is serious and this study does not
consider it refuted; it declines to rest any historical claim on either
answer.

\textbf{What would change this judgment.} For 1~Peter: a demonstrable
dependence on a document securely datable after 70, or a demonstrated use of
1~Peter by a witness earlier than any plausible secretarial window, would
settle the question in one direction or the other. For 2~Peter: a securely
dated witness quoting it before the end of the second century, or a
demonstration that Jude depends on 2~Peter rather than the reverse, would
substantially strengthen the case for its authenticity.
\end{studybox}
